\documentclass[10pt,a4paper]{amsart}
\usepackage[margin=19mm]{geometry}
\usepackage{amsmath,amssymb,mathtools}
\usepackage[scr=rsfs,cal=euler]{mathalfa}
\usepackage{xfrac}
\usepackage[unicode,hidelinks,bookmarks=false]{hyperref}
\newcommand{\ie}{i.e.\ }
\newcommand{\cf}{cf.\ }
\newcommand{\resp}{respectively\ }
\newcommand{\Subsection}{\subsection}
\providecommand{\nobreakdash}{\nobreak\hbox{-}}
\setlength{\emergencystretch}{2em}
\multlinegap0pt
\usepackage{amssymb}
\usepackage{mathtools}
\usepackage{tikz}
\newtheorem{lemma}{Lemma}
\title{Partial sums of random multiplicative functions \\with supercritical divisor twists}
\author{Jad Hamdan}
\date{Source: arXiv:2604.05563v1 (CC BY 4.0)}
\begin{document}
\maketitle

\noindent\emph{Source and licence.} Partial sums of random multiplicative functions \\with supercritical divisor twists. Version arXiv:2604.05563v1 (CC BY 4.0), distributed by arXiv under the Creative Commons Attribution 4.0 International licence, http://creativecommons.org/licenses/by/4.0/, as recorded in the arXiv licence metadata for this exact version and shown on the abstract page https://arxiv.org/abs/2604.05563v1. Full licence notice: \texttt{licenses/arxiv-2604.05563-CC-BY-4.0.txt}.\par\smallskip

\noindent\textit{Editorial scope.} Counted unit: the article's lemma lem:LDestimates (source0.tex lines 898--907). The article's complete original proof of it is delivered with it (source0.tex lines 908--967, one \texttt{proof} environment of 60 lines). No mathematics is written, repaired or re-derived: the statement and the proof are the article's own lines, in the article's own notation, and no retained line is substituted or re-flowed.  Retained uncounted as the source-local context the proof needs: lines 842--853.  Nothing was omitted from any retained span, so the package carries no omission record.  No retained line contains a citation, so the wrapper carries no bibliography and no bibliography entry.  No retained line contains a figure environment, an image-inclusion command or drawing code, so no image is delivered and the package declares no asset dependency.  Editorial formatting only, in addition: an AMS \texttt{amsart} preamble that restates the article's own declarations and macros used by the retained lines, namely the environments \texttt{lemma} on separate counters, together with the standard packages those lines need.  The retained environments print in this document as Lemma 1 in order of appearance, whereas the article prints the same items with the numbering of its own full document.  The complete native source of the article as distributed in the arXiv e-print for this exact version is bundled unchanged as \texttt{sources/197962/source0.tex}.
\begin{lemma}[Prime number theorem estimates]\label{lem:PNT}
Uniformly in $b>a>2$, there exists a $c>0$ for which
    \begin{equation}
        \sum_{a<p\leq b} \frac{1}{p}= \log\log b-\log\log a+O\big(e^{-c\sqrt{\log a}}\big),
    \end{equation}
and uniformly in $\sigma\in (0,1/2)$, 
\begin{equation}
    \sum_{a<p\leq b} \frac{1}{p^{2\sigma}} = \mathrm{Ei}\big((1-2\sigma)\log b\big)-\mathrm{Ei}\big((1-2\sigma)\log a\big)+O\big(e^{-c\sqrt{\log a}}\big)
\end{equation}
where $\mathrm{Ei}(x):=\int_{-\infty}^x(e^{s}/s)\mathrm{d}s$.
We also have that $\sum_{a\leq p\leq b}\frac{(\log p)^m}{p}=O\big((\log b)^m\big)$ for any $1\leq a\leq b$.
\end{lemma}

\begin{lemma}[Large deviation estimates]\label{lem:LDestimates} Let $C>0$ be an arbitrary constant. Then uniformly in all large $x$, $0\leq k\leq \log\log\log x$, $\sigma=1/2-k/\log x$, $1\leq j\leq \log\log x-k+1$, $|\alpha|\leq C$, and ${|v|\leq Cj}$,
\begin{equation}\label{eq:density}
        \mathbb{E}\big\{e^{\alpha S_j(\sigma)}\big\}\ll e^{\frac{\alpha^2}{4}j},\quad \mathrm{and} \quad \mathbb{P}\big(S_{j}(\sigma)\in [v,v+1)\big) \ll \frac{e^{-v^2/j}}{\sqrt{j}}.
\end{equation}
Furthermore, the same bounds hold under the measures $\mathbb{Q}=\mathbb{Q}_{h_1,h_2,\sigma,\beta}$ defined through
\[
    \frac{\mathrm{d}\mathbb{Q}}{\mathrm{d}\mathbb{P}}=\frac{\exp\big(\beta(S_j(\sigma+ih_1)-S_j(\sigma+ih_2))\big)}{\mathbb{E}\big\{\!\exp\big(\beta(S_j(\sigma+ih_1)-S_j(\sigma+ih_2))\big)\!\big\}}
\]
uniformly in $0\leq \beta\leq C|h_1-h_2|^{-1}e^{-j}$ and $h_1,h_2\in [-e^{-j-1},e^{-j-1}]$.
\end{lemma}

\begin{proof}
Since $|\alpha X_p^{(\sigma)}(0)|$ is bounded uniformly in $p\geq 2$ by our assumption on $\alpha$,  we can write
\[
    \log \mathbb{E}\big\{e^{\alpha S_j(\sigma)}\big\}=\sum_{p\leq \exp(e^j)} \log \mathbb{E}\big\{e^{\alpha X_p^{(\sigma)}(0)}\big\}=\sum_{p\leq \exp(e^j)} \log \Big(1+\frac{\alpha^2}{4p^{2\sigma}}+O\big(\alpha^3 p^{-3\sigma}\big)\Big)
\]
by Taylor expanding the exponential. Using Lemma \ref{lem:PNT} and the expansions $\log(1+x)=x+O(x^2)$ and $\mathrm{Ei}(x)=\gamma+\ln|x|+x+O(x^2)$,  
\begin{equation}\label{eq:MGFinc}
     \log \mathbb{E}\big\{e^{\alpha S_j(\sigma)}\big\}=\frac{\alpha^2}{4}\sum_{p\leq \exp(e^j)}\frac{1}{p^{2\sigma}}+O(\alpha^6p^{-3\sigma})\leq \frac{\alpha^2}{4}j+C',
\end{equation}
for some constant $C'$ (depending on $C$). It follows that $\mathbb{E}\{e^{\alpha S_j(\sigma)}\}\ll e^{(\alpha^2/4)j}$ for any $|\alpha|\leq C$. 

To estimate the probability in \eqref{eq:density}, we rewrite it as
\begin{equation}\label{eq:LD1}
    \mathbb{E}\big\{e^{\alpha S_j(\sigma)}\big\} \mathbb{E}_{\tilde{\mathbb{\mathbb{P}}}}\big\{e^{-\alpha S_j(\sigma)}\mathbf{1}\big(S_{j}(\sigma)\in [v,v+1)\big)\big\} \leq \mathbb{E}\big\{e^{\alpha S_j(\sigma)}e^{-\alpha v+\mathbf{1}(v<0)}\big\}\tilde{\mathbb{\mathbb{P}}}\big(S_{j}(\sigma)\in [v,v+1)\big)\big\}
\end{equation}
for $\alpha=2v/j$, where $\tilde{\mathbb{\mathbb{P}}}$ is the measure given by $\mathrm{d}\tilde{\mathbb{\mathbb{P}}}/{\mathrm{d}\mathbb{P}}={e^{\alpha S_{j}(\sigma)}}/{\mathbb{E}\{e^{\alpha S_{j}(\sigma)}}\}$. The expectation is $\ll e^{-v^2/j}$ by the earlier estimate $\mathbb{E}\{e^{\alpha S_j(\sigma)}\}\ll e^{(\alpha^2/4)j}$. To estimate the remaining probability, note that deterministically,
\[  
    \mathbb{E}_{\tilde{\mathbb{P}}}\Big\{\Big(S_j(\sigma)-\mathbb{E}_{\tilde{\mathbb{P}}}\big\{S_j(\sigma)\big\}\Big)^3\Big\}\leq C'\sum_{p} p^{-3\sigma} <\infty
\]
for some constant $C'>0$. Furthermore, there exists a constant $C_0$ depending only on $C$ such that 
\[  
    \frac{e^{\alpha X_p^{(\sigma)}(0)}}{\mathbb{E}\big\{e^{\alpha X_p^{(\sigma)}(0)}\big\}}=1+\alpha X_p^{(\sigma)}(0)+\frac{\alpha^2}{2}\Big(X_p^{(\sigma)}(0)^2-\mathbb{E}\big\{X_p^{(\sigma)}(0)^2\big\}\Big)+O(\alpha^3p^{-3\sigma}\big)
\]
uniformly for $p>C_0$, and we may assume that $C_0<\exp(e^j)$ without loss of generality (the desired bound otherwise follows by bounding the probability in \eqref{eq:LD1} by $1$). It follows that
\begin{align*}
    \mathrm{Var}_{\tilde{\mathbb{P}}}\big(S_j(\sigma)\big) &= O_C(1)+ \sum_{C_0<p\leq \exp(e^j)} \mathrm{Var}_{\tilde{\mathbb{P}}}\big(X_p^{(\sigma)}(0)\big) \\
    &= O_C(1)+\sum_{C_0<p\leq \exp(e^j)}\mathbb{E}\bigg\{\frac{e^{\alpha X_p^{(\sigma)}(0)}}{\mathbb{E}\big\{e^{\alpha X_p^{(\sigma)}(0)}\big\}}\Big(X_p^{(\sigma)}(0)-\mathbb{E}\big\{X_p^{(\sigma)}(0)\big\}\Big)^2\bigg\}\\
    &= O_C\big(1+\sum_p p^{-4\sigma}\big)+\mathrm{Var}_{\mathbb{P}}\big(S_j(0)\big) \asymp \mathrm{Var}_{\mathbb{P}}\big(S_j(0)\big).
\end{align*}
Using a standard Berry-Esseen bound and Lemma \ref{lem:PNT}, we conclude that
\[
    \tilde{\mathbb{\mathbb{P}}}\big(S_{j}(\sigma)\in [v,v+\Delta^{-1})\big)\ll \mathrm{Var}_{{\mathbb{P}}}\big(S_j(\sigma)\big)^{-1/2}\ll j^{-1/2}.
\]
The claim for $\mathbb{Q}$ follows from the same argument, provided one is armed with mean and variance estimates for $S_j(\sigma)$ under $\mathbb{Q}$, as well as a variance estimate under $\tilde{\mathbb{Q}}$ where $\mathrm{d}\tilde{\mathbb{\mathbb{Q}}}/{\mathrm{d}\mathbb{Q}}={e^{\alpha S_{j}(\sigma)}}/{\mathbb{E}_{\mathbb{Q}}\big\{e^{\alpha S_{j}(\sigma)}}\big\}$. We compute these directly, using the expansion
\[
    \frac{e^{\beta (X_p^{(\sigma)}(h_1)-X_p^{(\sigma)}(h_2))}}{\mathbb{E}\big\{e^{\beta (X_p^{(\sigma)}(h_1)-X_p^{(\sigma)}(h_2))}\big\}}=1+\beta \big(X_p^{(\sigma)}(h_1)-X_p^{(\sigma)}(h_2)\big)+O\Big(\beta^2 \big(X_p^{(\sigma)}(h_1)-X_p^{(\sigma)}(h_2)\big)^2\Big),
\]
for $p\leq \exp(e^j)$ large enough, and noting that the error term is
\[
   \ll  \frac{\beta^2|p^{-ih_1}-p^{-ih_2}|^2}{p^{2\sigma}} \leq \frac{\beta^2|h_1-h_2|^2(\log p)^2}{p^{2\sigma}} \ll \frac{1}{p^{2\sigma}}
\]
uniformly for $p\leq \exp(e^j)$ by our assumptions on $\beta$ and $|h_1-h_2|$. We therefore have

\begin{align}\label{eq:meanshift}
    \mathbb{E}_{\mathbb{Q}}\big\{S_j(\sigma)\big\} 
    &= O\Big(\sum_{p\leq \exp(e^j)} p^{-3\sigma}\Big)+\beta\sum_{p\leq \exp(e^j)} \mathbb{E}\big\{X_p^{(\sigma)}(0)\big(X_p^{(\sigma)}(h_1)-X_p^{(\sigma)}(h_2)\big)\big\}\notag\\
    &=O(1)+\frac{\beta}{2}\sum_{p\leq \exp(e^j)}
    \frac{\cos(h_1\log p)-\cos(h_2\log p)}{p^{2\sigma}}  \ll |h_1-h_2|\sum_{p\leq \exp(e^j)} \frac{\log p}{p} +1,
\end{align}
which is $O(1)$ by Lemma \ref{lem:PNT} since $|h_1-h_2|\leq e^{-j}$. Similarly, we find that
\begin{align}
    \mathrm{Var}_{\mathbb{Q}}\big(S_j(\sigma)\big)
    &= \mathrm{Var}_{\mathbb{P}}\big(S_j(\sigma)\big)
    + O\Big(\beta\,|h_1-h_2|\sum_{p} p^{-3\sigma}\log p\Big)
    + O\Big(\beta^2\sum_{p} p^{-4\sigma}\Big) \notag\\
    &= \mathrm{Var}_{\mathbb{P}}\big(S_j(\sigma)\big) + O(1),\label{eq:Qvar}
\end{align}
and $\mathrm{Var}_{\tilde{\mathbb{Q}}}\big(S_j(\sigma)\big)\asymp \mathrm{Var}_{\mathbb{Q}}\big(S_j(\sigma)\big)$.
\qedhere
\end{proof}

\end{document}
